% 格物（gewu）最小课题 Stage 5 成稿 —— 合成数据演示
% 生成自 output/min_case/model/summary.json 与 output/min_case/agents/gewu-improver/handoff_4_improver.json
% 编译：xelatex main.tex 或 tectonic main.tex（ctexart 需 UTF8 引擎）
\documentclass[11pt,a4paper,UTF8]{ctexart}
\usepackage[margin=2.5cm]{geometry}
\usepackage{booktabs}
\usepackage{amsmath}
\usepackage{hyperref}
\hypersetup{colorlinks=true,linkcolor=black,citecolor=blue,urlcolor=blue}
\geometry{margin=2.5cm}

\title{多示例学习辅助的 LIBS 大颗粒原煤热值预测\\[2pt]\large 合成数据上的最小链路验证与三臂交叉验证}
\author{格物科学智能体 (Gewu) $\cdot$ 最小课题合成演示}
\date{2026-11-10}

\begin{document}
\maketitle

\section{摘要}
\noindent\textbf{本文的实验对象为合成光谱数据，用于验证「检索—复现—改进—成稿—投稿—沉淀」七阶段链路可自主跑通；其绝对误差量级不代表真实煤样的精度结论。}
在 70 个 bag（共 611 条实例光谱）的 5 折 bag 级交叉验证下，袋均谱 PLSR 基线 RMSE 为
$150.44 \pm 13.15$ kcal/kg；仅做逐实例定标并池化（MIL 臂）为 $122.79 \pm 12.24$ kcal/kg；
进一步叠加强度上尾残差修正（本方法，hybrid 臂）为 $73.03 \pm 10.17$ kcal/kg，
相对基线降低 $51.46\%$。逐折配对 $t$ 检验给出 $t(4)=-22.48$、$p=2.3\times 10^{-5}$，
95\% 置信区间为 $[-86.97, -67.85]$ kcal/kg，五折方向一致为改进；
Wilcoxon 符号秩检验 $p=0.0625$（$n=5$ 时双侧最小可达 $p$ 值，本身不显著），
两种检验结论一并报告，不以此宣称强统计显著性。
方法学上的关键观察是：内层交叉验证选出的 LSE 温度参数 $\beta$ 有 4 折收敛到 $0.25$、
1 折收敛到 $0$，故增益不应被归因于高阶池化，而应归因于逐实例定标对上尾颗粒异质性信息的回收。

\section{引言}
激光诱导击穿光谱（LIBS）用于煤质快速分析时，样品以「大颗粒堆积」形式被烧蚀：
一个煤样对应多条光谱（多实例），而标签只在\textbf{样}级可得，单条光谱没有独立标签。
这正是多示例学习（Multiple Instance Learning, MIL）的标准设定\cite{ilse2018mil}。

在煤质 LIBS 定量领域，主流做法是把一个样品的多条光谱先平均成一条「袋均谱」，
再做偏最小二乘回归（PLSR）\cite{li2014pls,yan2023coal}。
袋均操作在数学上等价于对实例取期望：$\overline{\mathbf{x}}_j=\frac{1}{n_j}\sum_{i\in j}\mathbf{x}_i$。
当颗粒尺度异质性 $h$ 只在样品级起作用时，$\overline{a_i}\to 1$ 使袋均谱几乎把 $h$ 完全抹掉；
少数粗颗粒主导信号，平均后这部分信息被稀释，成为 PLSR 的不可约残差。

本文的贡献不在提出新网络，而是给出一条\textbf{可被第三方独立重跑}的最小链路：
数据可合成、模型只依赖 numpy 与 scikit-learn、每个量化断言都能回溯到
JSON Pointer（见本文第 5 节表 1 对应的 claims 清单）。
具体地，本文把「成分」与「颗粒尺度」在合成生成机制中显式分离，
并对照三个共享同一折划分与同一批特征的臂，以评估逐实例定标与残差修正各自的贡献。

\section{方法}
\subsection{合成数据生成机制}
合成光谱按式 (1) 生成，把成分向量 $\mathbf{c}_i$（逐实例可见）与粒度因子 $a_i$ 相乘：
\begin{equation}
  \mathbf{x}_i = a_i \cdot \bigl[\, \mathbf{b} + \sum_k c_k \mathbf{p}_k \,\bigr] \cdot (1+\varepsilon_i) + \mathbf{n}_i,
  \qquad a_i = 1 + h\,z_i,
\end{equation}
其中 $h$ 为样品级颗粒异质性，\textbf{仅 bag 级可见}（同一 bag 内实例共享 $h$）；
$z_i$ 为逐实例尺度扰动，$\varepsilon_i$ 与 $\mathbf{n}_i$ 为加性噪声。
标签按式 (2) 给出：
\begin{equation}
  y = 3000 + 1200\,c_0 + 900\,c_1 + 600\,c_2 + 1400\,h + \mathcal{N}(0, 20).
\end{equation}
关键性质：袋均运算使 $\overline{a_i}\to 1$，从而抹掉 $h$；
而上尾统计量（$\max_i a_i$ 等）保留 $h$，这正是残差修正的信息来源。

\subsection{三个对照臂}
三臂共享同一 5 折 bag 级划分与同一批原始特征，仅改变「定标方式 / 池化 / 残差建模」：
\begin{itemize}
  \item \textbf{PLSR 基线}：袋均谱 $\to$ PLSR（主成分数 $n_c=6$）；
  \item \textbf{MIL 直接}：对每条实例谱做谱形定标 $\mathbf{x}_i/\overline{\mathbf{x}_i}$，
        经实例级 PLS 后由 LSE 池化聚合；LSE 温度参数 $\beta$ 由内层 4 折交叉验证选出；
  \item \textbf{本方法（hybrid）}：在 MIL 直接之上，另取强度池化统计量
       （$\max$、$P_{90}$、$\mathrm{std}$、极差、实例数，共 7 维）作为特征，
        用 Ridge 回归 PLSR 的折外残差。
\end{itemize}
残差特征刻意只取基线「看不到」的量，不放入与 MIL 输出共线的预测统计量，
否则会在 56 个训练 bag 上过拟合。

\section{实验设置}
\subsection{数据与折划分}
数据集含 70 个 bag、611 条实例光谱，按 bag 切分为 5 折（每折 14 个 bag）：
同一 bag 的全部实例只落入同一折，避免信息泄漏。
随机种子固定，脚本化重跑结果逐位一致。

\subsection{超参数选择口径}
$\beta$ 与 Ridge 的 $\alpha$ \textbf{只在内层 4 折的折外预测上选择}；
残差模型的输入特征同样取自折外预测，而非训练集自拟合，
以避免乐观偏差（optimistic bias）。
所有指标按「每折各自算 RMSE，再对 5 折取均值与总体标准差（$\mathrm{ddof}=0$）」的口径报告；
同时给出 70 个袋外预测合并后的 pooled 指标。两种口径不可混用。

\section{结果}
\begin{table}[h]
\centering
\caption{5 折交叉验证三臂汇总（RMSE 单位：kcal/kg；均值为折均值，标准差为总体标准差）}
\label{tab:cv}
\begin{tabular}{lcccc}
\toprule
臂 & 均值 & 总体标准差 & $R^2$ 均值 & 逐折 RMSE \\
\midrule
PLSR 基线 & 150.4354 & 13.1549 & 0.6106 & 169.6308 / 134.7206 / 158.1094 / 136.8353 / 152.8811 \\
MIL 直接 & 122.7867 & 12.2446 & 0.7407 & 141.3063 / 109.2403 / 118.5462 / 112.2981 / 132.5425 \\
本方法 & 73.0270 & 10.1693 & 0.9108 & 91.8709 / 69.1661 / 72.3988 / 61.1779 / 70.5212 \\
\bottomrule
\end{tabular}
\end{table}

本方法相对 PLSR 基线 RMSE 下降 $77.4084$ kcal/kg（$-51.46\%$）；
相对 MIL 直接臂再下降 $33.08\%$（MIL 直接相对基线为 $-18.38\%$）。
袋外合并（pooled）指标：PLSR RMSE $=151.0095$、$R^2=0.6309$；
MIL 直接 RMSE $=123.3957$、$R^2=0.7535$；本方法 RMSE $=73.7316$、$R^2=0.912$。

逐折配对 $t$ 检验（本方法减基线）：$t=-22.4798$，自由度 $4$，
双侧 $p=2.3\times 10^{-5}$，均值差 95\% 置信区间 $[-86.969, -67.8479]$ kcal/kg，
五折差值分别为 $-77.7599$、$-65.5545$、$-85.7106$、$-75.6574$、$-82.3599$，方向一致。
Wilcoxon 符号秩检验统计量 $0$、$p=0.0625$；
鉴于 $n=5$ 时双侧最小可达 $p$ 即为 $0.0625$，Wilcoxon 本身无法达到 $0.05$ 显著性，
本文把两种检验结论并列，结论以「方向一致 + 效应量 $51.46\%$ + $t$ 检验显著」合并表述。

\section{消融与讨论}
三臂对比唯一可控的变化是「是否做逐实例定标」与「是否叠加强度上尾残差」，
因此可做归因分解：去掉残差（MIL 直接相对基线）贡献 $-18.38\%$，
再加残差贡献额外的 $-33.08\%$。
这说明本合成数据上的增益主要来自两点：
(1) \textbf{逐实例定标}把尺度因子与成分解耦；
(2) \textbf{强度上尾统计量}回收了被袋均抹掉的粒度异质性 $h$。

\subsection{池化温度的诚实解读}
内层交叉验证选出的 $\beta$ 为 $[0, 0.25, 0.25, 0.25, 0.25]$。
$\beta=0$ 等价于均值池化，因此 5 折中 1 折明确选中了均值池化、
其余 4 折选中了一个小的非零温度。
\textbf{若 $\beta=0$ 占多数，应如实解读为「该数据上均值池化已足够，增益不来自高阶池化」，
而不是宣称 LSE 有效。}本文据此不把增益归因于池化算子本身。

\subsection{未覆盖的消融}
组件级消融（逐个去掉 $\max$、$P_{90}$、$\mathrm{std}$ 等残差特征）未做：
在 56 个训练 bag 上逐一关组件易过拟合，结论不可靠，故留作后续。
此外，本课题未在任何真实煤样上验证，需外部私有数据，不在本课题范围内。

\subsection{适用边界}
合成数据不含基体效应、自吸收与仪器漂移，上述误差量级不可外推到真实煤样。
本报告证明的是\textbf{机制与可复现链路}，真实精度以存档课题的实测指标为准。

\section{结论}
在零外部数据依赖的前提下，七阶段链路可从零跑出可复现的数值、
可回溯的断言与可编译的成稿：
(1) 三臂共享折划分的 5 折交叉验证给出本方法 RMSE $73.03\pm 10.17$ kcal/kg，
相对 PLSR 基线降低 $51.46\%$，五折方向一致；
(2) 配对 $t$ 检验 $t(4)=-22.48$、$p=2.3\times 10^{-5}$，
Wilcoxon $p=0.0625$ 一并报告，不宣称强统计显著性；
(3) 增益归因为逐实例定标与上尾残差，而非 LSE 池化（$\beta$ 有 1 折为 0）；
(4) 每个量化断言均带 $\{$文件, JSON Pointer$\}$ 溯源，可由第三方独立核对。
本链路的终点是「成果可复现性」的第三方独立可执行证据，而非真实煤样的精度声明。

\begin{thebibliography}{9}
\bibitem{li2014pls} Y. Li, Z. Wang, T. Li, et al. A model combining spectrum standardization and dominant factor based partial least square method for carbon analysis in coal by laser-induced breakdown spectroscopy. \textit{Spectrochimica Acta Part B: Spectroscopy}, 100:74--80, 2014. DOI: 10.1016/j.sab.2014.06.017.
\bibitem{ilse2018mil} M. Ilse, J. Tomczak, M. Welling. Attention-based Deep Multiple Instance Learning. In \textit{Proc. ICML}, 2018. arXiv:1802.04712.
\bibitem{yan2023coal} Y. Yan, et al. Rapid Classification and Quantification of Coal by Using Laser-Induced Breakdown Spectroscopy and Machine Learning. \textit{Applied Sciences}, 13(14):8158, 2023. DOI: 10.3390/app13148158.
\end{thebibliography}

\end{document}